% =========================================================
% Section: Secant and Tangent Lines
% =========================================================

\begin{tcolorbox}[colback=gray!8, colframe=gray!40, title=\textbf{Secant and Tangent --- toolkit}]
\textbf{Concept family:} The geometric objects underlying the derivative: secant lines and tangent lines.

\medskip
\textbf{Atomic items in this section:}
\begin{itemize}
  \item Definition: Secant Line
  \item Definition: Tangent Line
\end{itemize}
\end{tcolorbox}

% ---------------------------------------------------------

\begin{tcolorbox}[colback=blue!4, colframe=blue!30]
\begin{definition}[Secant Line]
\label{def:secant-line}
Let $f : A \to \mathbb{R}$, and let $x_1, x_2 \in A$ with $x_1 \neq x_2$. The \emph{secant line} determined by the points $(x_1, f(x_1))$ and $(x_2, f(x_2))$ is the unique line passing through these two points. Its slope is
\[
m_{\mathrm{sec}}(x_1, x_2) := \frac{f(x_2) - f(x_1)}{x_2 - x_1}.
\]
\end{definition}
\end{tcolorbox}

\begin{remark*}[Standard quantified statement]
\[
\forall x_1, x_2 \in A\;(x_1 \neq x_2 \Rightarrow \operatorname{SecantLine}(f, x_1, x_2)
\text{ has slope } \tfrac{f(x_2)-f(x_1)}{x_2-x_1}).
\]
\end{remark*}

\begin{remark*}[Interpretation]
With increments $\Delta x := x_2 - x_1$ and $\Delta y := f(x_2) - f(x_1)$, the secant slope is $\Delta y / \Delta x$ --- the average rate of change of $f$ over the interval $[x_1, x_2]$. This is an exact algebraic identity; no limit is involved. As $x_2 \to x_1$, if the secant slope converges, its limit is the derivative $f'(x_1)$ --- the instantaneous rate of change. Geometrically, the secant line rotates about the fixed point $(x_1, f(x_1))$ as $x_2$ approaches $x_1$, and the limit position is the tangent line.
\end{remark*}

\NoLocalDependencies

% ---------------------------------------------------------

\begin{tcolorbox}[colback=blue!4, colframe=blue!30]
\begin{definition}[Tangent Line]
\label{def:tangent-line}
Let $f : A \to \mathbb{R}$, let $c \in A$, and suppose that $f$ is differentiable at $c$. The \emph{tangent line} to the graph of $f$ at the point $(c, f(c))$ is the line through $(c, f(c))$ with slope $f'(c)$. Its equation is
\[
y = f(c) + f'(c)(x - c).
\]
\end{definition}
\end{tcolorbox}

\begin{remark*}[Standard quantified statement]
\[
\operatorname{DifferentiableAt}(f, c) \;\Rightarrow\;
\operatorname{LinearApproximationAt}(f, c) = f(c) + f'(c)(x - c).
\]
\end{remark*}

\begin{remark*}[Interpretation]
The tangent line is the limit of the family of secant lines through $(c, f(c))$ and $(x, f(x))$ as $x \to c$. More precisely: the slope $m_{\mathrm{sec}}(c, x) = (f(x)-f(c))/(x-c)$ converges to $f'(c)$, so the tangent line is the line whose slope is that limit. It is also the unique linear function $\ell(x) = f(c) + f'(c)(x-c)$ satisfying $\ell(c) = f(c)$ and $\ell'(c) = f'(c)$ --- the best first-order approximation to $f$ near $c$. The error $f(x) - \ell(x) = o(x-c)$ as $x \to c$, meaning it vanishes faster than the displacement.
\end{remark*}

\begin{remark*}[Dependencies]
\begin{itemize}
  \item \hyperref[def:secant-line]{Definition: Secant Line}
  \item \hyperref[def:derivative-at-a-point]{Definition: Derivative at a Point}
\end{itemize}
\end{remark*}

% Figure: secant and tangent lines --- extracted to figure file
\input{volume-iii/analysis/differentiation/notes/secant-tangent/figure1}
